% 张三慧《学习辅导与习题解答》PDF S 独立选答，卷1第2章；未取得V1主教材独立交叉核验。
% 仅录“思考题选答”中实际明确列出的原编号；其余不推断，exam_ready=false。
\section*{第2章\quad 运动与力：辅导书已列思考题独立解析}
\paragraph{思考题 2.6}
木板自由下落的瞬时随动参考系加速度近似 \(g\) 向下，重力与惯性力抵消，摆球只受张力，若原有非零相对切向速度，可绕悬点作等速圆周运动（忽略阻力和板侧扰动）；若割断时刚在极端位置，相对切向速度为零，则在理想自由落体参考系里摆球可相对木板保持静止。该结论依赖整个装置自由落体及小摆长近似。

\paragraph{思考题 2.7}
汽车向左转弯时，其随动车厢系中存在指向右侧的离心惯性加速度，空气等效重力 \(\boldsymbol g_{\rm eff}\) 的水平分量向右。空气压力沿等效重力方向升高，氦气球受浮力沿\(-\boldsymbol g_{\rm eff}\)，水平分量因此\emph{向左，即向弯道内侧}。这与普通较重悬挂物的外甩方向相反。

\paragraph{思考题 2.11}
若卫星刚体不同部分离地心距离稍有差异，万有引力随 \(r^{-2}\) 变化会产生\emph{重力梯度力矩}，其量级在近圆轨道为 \(\boldsymbol\tau=3\mu r^{-3}\hat{\boldsymbol r}\times(\mathbf I\hat{\boldsymbol r})\)（取符号约定）。长轴指向地球的姿态在适当惯量排序时成为小扰动可回复的稳定取向。只说“长杆引力较大”不充分；需核对姿态刚体的主惯量与动态稳定条件。
